\chapter{Case Study: Developing a Model Checker}
\label{ch:model-checker-case-study}

A code cache can contain machine code that was correct when compiled and unsafe when called.
Suppose an optimization depends on a class hierarchy remaining unchanged.
The runtime publishes an optimized entry, later invalidates the hierarchy assumption, and eventually removes or replaces the entry.
An unchecked call between invalidation and removal can enter stale code.
Testing a successful compile-and-call sequence does not exercise that ordering.
The missing test is a sequence of protocol transitions, rather than a different arithmetic input to the generated function.

A small \term{finite-state model checker} makes these sequences explicit and searches them systematically.
The software developed here accepts a finite labeled graph, checks whether any initial state can reach a bad state, and returns a shortest failing execution when one exists.
It also evaluates three temporal operators over sets of states.
The complete C11 lab is in \code{example-code/model-checker/}; its input files, command-line interface, independent test oracles, and exercise answers accompany the implementation.
The checker uses only the C standard library and has a declared capacity of 64 states.
That bound keeps the representation inspectable without pretending that arbitrary compiler behavior has become finite or small.

\section{Choose a property and an abstraction boundary}
\label{sec:ch48-model}

Our property is precise: no execution starting in the interpreted state may invoke optimized code after its assumption has been invalidated.
We represent one cache entry, one relevant assumption, and one call at a time.
Compilation, invalidation, invocation, and retirement are atomic model steps.
We do not represent instructions inside compiled code, the data reconstructed by deoptimization, or concurrent reads of a publication pointer.
Those choices determine what a successful check means.

Table~\ref{tab:ch48-states} assigns stable IDs to the six model states.
The proposition mask belongs to a state; it is not the state's numeric ID.
Bit zero means \code{bad}, and bit one means \code{completed}.
The second proposition is available to callers even though the CLI's safety query uses only the first.

\begin{table}[htbp]
  \caption{The cache protocol records validity and invocation outcomes separately.}
  \label{tab:ch48-states}
  \centering
  \begin{tabular}{clc}
    \toprule
    ID & Meaning & Proposition mask \\
    \midrule
    0 & Interpreted execution; no optimized entry & 0 \\
    1 & Optimized entry published; assumption valid & 0 \\
    2 & Assumption invalidated; entry still present & 0 \\
    3 & Unchecked invocation entered stale code & 1 \\
    4 & Valid optimized call completed & 2 \\
    5 & Guard rejected entry; deoptimized call completed & 2 \\
    \bottomrule
  \end{tabular}
\end{table}

Figure~\ref{fig:ch48-cache} contains the whole transition relation.
State 0 is the only initial state.
From state 1 a valid call can complete, or the assumption can be invalidated first.
From state 2 the buggy protocol permits both a guarded fallback and an unchecked invocation.
Nondeterminism represents the possibility that different call paths enforce the guard differently.
The corrected protocol removes the unchecked edge $2\to3$: every call in state 2 must take the guarded deoptimization path.
State 3 remains a named state in both files, so we can distinguish unreachable errors from errors deleted from the vocabulary.

\begin{figure}[htbp]
  \centering
  \begin{tikzpicture}
  \node[basic block] (s0) at (0,0) {0: interpret};
  \node[basic block] (s1) at (3,0) {1: valid};
  \node[basic block] (s2) at (6,0) {2: stale};
  \node[basic block] (s3) at (9,0) {3: bad};
  \node[basic block,align=center] (s4) at (3,-2) {4: valid call\\completed};
  \node[basic block,align=center] (s5) at (6,-2) {5: deoptimized\\call completed};
  \draw[cfg edge] (s0) -- node[above,font=\scriptsize] {publish} (s1);
  \draw[cfg edge] (s1) -- node[above,font=\scriptsize] {invalidate} (s2);
  \draw[back edge] (s2) -- node[above,font=\scriptsize] {unchecked call} (s3);
  \draw[cfg edge] (s1) -- node[left,font=\scriptsize] {valid call} (s4);
  \draw[cfg edge] (s2) -- node[right,font=\scriptsize] {guard / deoptimize} (s5);
  \draw[cfg edge] (s4.west) to[bend left=20] node[left,font=\scriptsize] {retire} (s0.south);
  \draw[cfg edge] (s5.south) -- ++(0,-0.65) -| (s0.south);
  \draw[cfg edge] (s3) edge[loop right] node[right,font=\scriptsize] {implicit} (s3);
  \draw[cfg edge] (-1.7,0) -- (s0);
\end{tikzpicture}

  \caption{The dashed edge admits the stale-code bug. Removing that edge makes the guarded invocation mandatory; the error state's loop is supplied by the dead-end convention.}
  \label{fig:ch48-cache}
\end{figure}

This graph is a model of a compiler runtime protocol, not a proof of a whole compiler.
An implementation can violate a modeled atomic step, omit an invalidation event, or reconstruct the wrong frame while following the graph's legal edges.
To transfer a safety result to a real runtime, we need an abstraction map from concrete executions to model executions and a reason that concrete errors are represented by bad model states.
An abstraction that omits possible concrete transitions can falsely report safety.
An abstraction that adds transitions may report an infeasible counterexample, which must then be investigated or refined.
The graph checker establishes a fact about the supplied graph; the modeling argument connects that fact to software.

\section{Make input and ownership part of the design}
\label{sec:ch48-interface}

Figure~\ref{fig:ch48-pipeline} separates parsing, validation, query evaluation, and result consumption.
This separation is useful even in a small program.
A malformed file must not become a partially constructed graph that happens to have no reachable error.
A budget failure must not become an empty satisfaction set interpreted as a completed analysis.
The output status therefore controls whether an output value has mathematical meaning.

\begin{figure}[htbp]
  \centering
  \begin{tikzpicture}[node distance=8mm and 10mm]
  \node[basic block,align=center] (input) {decimal graph\\caller-owned stream};
  \node[basic block,align=center,right=of input] (read) {parse + validate\\temporary graph};
  \node[basic block,align=center,right=of read] (graph) {published graph\\immutable input};
  \node[basic block,align=center,below=of graph] (kernel) {BFS / CTL kernel\\explicit budget};
  \node[basic block,align=center,left=of kernel] (result) {status + result\\trace or state set};
  \node[basic block,align=center,left=of result] (replay) {caller checks status\\replay / compose};
  \draw[cfg edge] (input) -- (read);
  \draw[cfg edge] (read) -- node[above,font=\scriptsize] {success} (graph);
  \draw[cfg edge] (graph) -- (kernel);
  \draw[cfg edge] (kernel) -- (result);
  \draw[cfg edge] (result) -- (replay);
\end{tikzpicture}

  \caption{The parser publishes a graph only after validation. Query results carry a status that callers must inspect before replaying a trace or composing a state set.}
  \label{fig:ch48-pipeline}
\end{figure}

The graph file is a sequence of unsigned decimal integers separated by whitespace.
Its header contains $N$, the number of states; $K$, the number of initial states; and $M$, the number of edge records.
The next $K$ integers are distinct initial state IDs.
The next $N$ integers are proposition masks in state order.
The remaining $2M$ integers form source--destination pairs.
The parser then requires whitespace followed by end of file.
Listing~\ref{lst:ch48-model-file} is the actual failing model consumed by the tests, not a second transcription of it.

\begin{listing}[htbp]
  \inputminted{text}{example-code/model-checker/stale.graph}
  \caption{The failing model has six states, one initial state, and seven explicit edges.}
  \label{lst:ch48-model-file}
\end{listing}

IDs must lie in $[0,N)$; $1\leq N\leq64$, $1\leq K\leq N$, and $0\leq M\leq4096$.
Duplicate edges are harmless and coalesce into one bit.
Duplicate initial IDs are rejected, so the declared initial-state count is meaningful.
Signs, hexadecimal prefixes, nonnumeric suffixes, truncated records, and trailing tokens are invalid.
Masks may use every bit of an unsigned 64-bit integer.
Before accumulating decimal digit $d$ into value $v$, the parser checks
\[
  v \leq \left\lfloor\frac{\mathrm{UINT64\_MAX}-d}{10}\right\rfloor.
\]
Only then does it evaluate $10v+d$.
An integer that cannot be represented is invalid; a representable header exceeding the supported state or edge capacity produces a resource-limit result.
The parser can stop at that header and need not consume the remaining stream.

The graph contains an initial-state bitset, one proposition mask per state, and one successor bitset per state.
Listing~\ref{lst:ch48-api} shows the public types and query signatures.
A proposition mask and a state set are both represented by \code{uint64\_t}, but they index different things.
The function \code{mc\_atoms(g, mask)} converts the former to the latter: it returns states having \emph{any} proposition bit selected by \code{mask}.
Passing zero selects no states.
Explicit conversion avoids accidentally interpreting proposition bit zero as state zero.

\begin{listing}[htbp]
  \inputminted[firstline=9,lastline=32]{c}{example-code/model-checker/checker.h}
  \caption{Caller-owned fixed-size objects make graph and trace lifetimes explicit.}
  \label{lst:ch48-api}
\end{listing}

There is no heap allocation and no hidden graph handle.
The caller owns the graph, trace, and result bitset; input and output objects must be distinct and nonoverlapping.
No function retains a pointer after returning, and no destruction function is needed.
The parser clears its output, builds a temporary graph, and copies that graph to the output only on success.
Queries validate graph dimensions, the nonempty initial set, and successor masks before indexing arrays.
Validation assumes live C objects: it cannot determine whether an arbitrary pointer denotes an allocation of the advertised type.
The graph is immutable during a query, which also makes independent concurrent queries possible when their output objects are separate.

The status contract has five values.
\code{MC\_OK} means that parsing or a temporal-set operation completed.
\code{MC\_SAFE} and \code{MC\_UNSAFE} belong to safety reachability.
\code{MC\_INVALID} reports invalid input, and \code{MC\_LIMIT} reports capacity or query-budget exhaustion.
Every non-unsafe safety result leaves an empty trace.
Every unsuccessful CTL operation leaves a zero output set, but that zero is not a semantic answer until the caller has checked the status.
The CLI maps safe, unsafe, invalid, and resource-limit outcomes to exit codes 0, 1, 2, and 3 respectively.

\section{Define what a dead end means}
\label{sec:ch48-totalization}

A finite graph may contain a state with no successors, but the temporal operators below quantify over infinite paths.
We resolve that mismatch with \term{totalization}: a state with no explicit successor has an implicit edge to itself.
Existing edges are unchanged.
The library computes this effective successor set on demand instead of mutating the graph.
The chosen interpretation says that after execution stops, its final observation persists.

This convention affects temporal formulas.
At a dead end satisfying $p$, both $EX\,p$ and $EG\,p$ hold because the repeated final state is an infinite witness.
At a dead end not satisfying $p$, neither holds.
Safety reachability is unchanged: the extra self-loop cannot make a different state reachable.
If termination or deadlock is itself an error, the model exporter must mark the corresponding states before totalization.
A formula observing only the proposition labels cannot distinguish a supplied self-loop from an implicit one.
The checker imposes no fairness assumption; a permitted cycle may repeat forever even if another transition is continuously available.

\section{Return a shortest safety counterexample}
\label{sec:ch48-bfs}

Let $I$ be the initial-state set and $B$ the states whose proposition masks intersect the requested bad mask.
The safety question is whether any state in $B$ is reachable from any state in $I$.
Algorithm~\ref{alg:ch48-bfs} uses a multi-source breadth-first search.
Every initial state enters the queue before exploration starts, so all of them have distance zero.
Each newly discovered state receives exactly one parent, and each root is its own parent.

\begin{algorithm}[htbp]
  \caption{Budgeted safety search with a replayable shortest trace.}
  \label{alg:ch48-bfs}
  \begin{pseudocode}
    \Require Valid graph, nonempty initial set $I$, bad set $B$, dequeue budget $L$
    \Ensure Safe, unsafe with a shortest trace, or resource limit
    \State $Q\gets$ initial states in increasing order; $D\gets I$; $v\gets0$
    \ForAll{$s\in I$}
      \State $\operatorname{parent}[s]\gets s$
    \EndFor
    \While{$Q$ is not empty}
      \If{$v=L$}
        \State \Return resource limit
      \EndIf
      \State $s\gets\operatorname{popFront}(Q)$; $v\gets v+1$
      \If{$s\in B$}
        \State Follow parents from $s$ to its root and reverse the sequence
        \State \Return unsafe with that sequence
      \EndIf
      \ForAll{effective successors $t$ of $s$, in increasing order}
        \If{$t\notin D$}
          \State $D\gets D\cup\{t\}$; $\operatorname{parent}[t]\gets s$
          \State Append $t$ to $Q$
        \EndIf
      \EndFor
    \EndWhile
    \State \Return safe
  \end{pseudocode}
\end{algorithm}

The critical invariant is marking on enqueue, not on dequeue.
Once a state is marked, no later edge appends it or changes its parent.
The queue therefore contains at most $N$ entries, and parent chains cannot acquire cycles except the intentional root self-parent.
Queue order is nondecreasing shortest distance from the initial set.
The first dequeued bad state consequently gives a shortest counterexample in number of transitions.
An initially bad state produces a one-state trace, representing zero transitions.
Increasing-ID iteration makes equal-length choices reproducible, without claiming a preferred semantic explanation among several shortest failures.

For the buggy cache model, the successive queue contents, retaining already examined entries, are
\[
  [0],\quad [0,1],\quad [0,1,2,4],\quad [0,1,2,4,3,5].
\]
Examining state 4 adds nothing because state 0 has already been discovered.
The next examination reaches state 3.
Its parent chain is $3,2,1,0$; reversing it returns $0,1,2,3$.
The trace has three edges: publish, invalidate, and unchecked invocation.
Notice that five dequeues were needed, because state 4 preceded state 3 in the queue.
A budget measures work, not the length of the eventual witness.

Listing~\ref{lst:ch48-bfs-core} includes the executed dequeue, error test, reconstruction, and discovery code.
Trace replay checks that its first state is initial, each consecutive pair is an edge in the effective graph, and its last state is bad.
The sequence contains state IDs rather than memory addresses or internal queue positions.
It can be saved in a regression test and resolved against a stable model export.
Because this graph has no transition labels or parallel-edge identities, replay identifies a state path, not a unique external action sequence in a richer system.

\begin{listing}[htbp]
  \inputminted[firstline=109,lastline=135]{c}{example-code/model-checker/checker.c}
  \caption{The BFS core publishes a trace only when it has reached a bad state.}
  \label{lst:ch48-bfs-core}
\end{listing}

The implementation scans all $N$ possible destinations for each examined state.
Its worst-case safety-query cost is therefore $O(N^2)$, not the $O(N+M)$ cost of an adjacency-list BFS.
The queue, parent array, and trace use $O(N)$ machine words; graph storage uses $N$ successor words and $N$ proposition words for $N\leq64$.
A sparse-graph implementation can replace successor scans with adjacency lists while preserving the queue and parent invariants.
This lab chooses a small bounded representation so that bounds checks and bit operations remain visible.

\section{Build temporal queries from state sets}
\label{sec:ch48-ctl}

Safety asks whether an error can happen at any future point.
Other useful questions ask whether an event can occur on the next step, whether a condition can be maintained until a goal is reached, or whether some execution can remain in a region forever.
The lab implements the existential next, until, and globally operators of computation tree logic, abbreviated CTL, using the standard finite-state semantics described by Baier and Katoen~\cite{baier-2008-model-checking}.
Each operand and result is a set of states.
The operators are evaluated over all states, including those unreachable from the initial set.

For state set $X$, define its existential predecessor set by
\[
  \operatorname{Pre}(X)
  = \{s\mid\text{some effective successor of }s\text{ belongs to }X\}.
\]
The operator $EX\,p$ holds where at least one immediate successor satisfies $p$.
Thus its satisfaction set is $\operatorname{Pre}(P)$ when $P$ is the satisfaction set of $p$.
The implementation needs one scan of the successor words.

The operator $E[p\ U\ q]$ requires a path with a finite position satisfying $q$, with $p$ holding at every earlier position.
The endpoint need not satisfy $p$, and a current state satisfying $q$ is already a witness of length zero.
Its least-fixpoint computation starts at $X_0=Q$ and iterates
\[
  X_{i+1}=Q\cup\bigl(P\cap\operatorname{Pre}(X_i)\bigr).
\]
Here $X_i$ contains states with witnesses of at most $i$ transitions.
The sequence grows monotonically, and every strictly changing round adds at least one state.
Starting from the empty set and applying the same transfer function is also correct, but spends an extra initial round installing $Q$.

The operator $EG\,p$ requires an infinite path on which $p$ always holds.
Its greatest-fixpoint computation starts at $Y_0=P$ and iterates
\[
  Y_{i+1}=P\cap\operatorname{Pre}(Y_i).
\]
Each round removes states that cannot stay in the current candidate region for another step.
At convergence, every surviving state has a surviving successor.
Repeatedly choosing such a successor produces an infinite path in $P$.
Conversely, an infinite path entirely in $P$ cannot lose its states during this elimination, because each state on the path retains a qualifying successor.
For a finite graph this also means that a qualifying path can reach a cycle entirely inside $P$.

Listing~\ref{lst:ch48-ctl-core} is the complete temporal-operation entry point.
The loop maintains a private approximation and writes the caller's result only after completion.
EU and EG require at most $N$ changing rounds and one equality check, so $N+1$ is a sufficient budget.
One predecessor operation costs $O(N)$ word operations in this representation, giving an $O(N^2)$ worst-case bound for either fixpoint and $O(N)$ for EX, including input validation.
These bounds assume a machine word holding all states; extending the capacity requires a multiword bitset or a different graph representation and a revised cost analysis.

\begin{listing}[htbp]
  \inputminted[firstline=146,lastline=168]{c}{example-code/model-checker/checker.c}
  \caption{One entry point implements EX and the ascending EU and descending EG iterations.}
  \label{lst:ch48-ctl-core}
\end{listing}

On the buggy graph, $E[\mathit{true}\ U\ \mathit{bad}]$ grows through
\[
\begin{split}
  \{3\}&\subset\{2,3\}\subset\{1,2,3\}\\
         &\subset\{0,1,2,3\}\subset\{0,1,2,3,4,5\}.
\end{split}
\]
A fifth sweep confirms equality with the final set.
Four sweeps produce the same bitset internally but do not establish convergence in this implementation, so a four-round budget returns \code{MC\_LIMIT} with no published answer.
In the repaired graph the starting set $\{3\}$ is already a fixpoint; one sweep confirms it.
The repaired model still has a bad state, but no initial state can reach it.

Boolean composition uses union, intersection, and complement relative to the universe of $N$ states.
For example, the model satisfies a state formula with set $F$ exactly when $I\subseteq F$.
Universal safety can be checked by testing
\[
  I\cap\llbracket E[\mathit{true}\ U\ \mathit{bad}]\rrbracket=\varnothing.
\]
In contrast, $EG\,\neg\mathit{bad}$ merely asserts the existence of an infinite safe path.
It holds at state 0 even in the buggy model: the cycle $0\to1\to4\to0$ never enters stale code.
Confusing existential survival with universal safety would therefore hide the bug.
EX, EU, EG, and Boolean operations supply a basis for CTL state formulas, but the lab exposes an operator API rather than a formula-language parser.
It does not implement general LTL checking, fairness constraints, or temporal witness extraction for EU and EG.

\section{Integrate the checker into a compiler workflow}
\label{sec:ch48-integration}

The CLI takes a graph path and an optional maximum number of BFS dequeues.
Its default budget is 64, sufficient to examine every state of any accepted graph.
After a completed safety query, it also prints three CTL state sets using an independent $N+1$-round budget.
From the repository root, the following commands build and run the failing example; the nonzero exit status is intentional.

\begin{verbatim}
gcc -std=c11 -Wall -Wextra -Wpedantic -Werror -O2 \
  example-code/model-checker/checker.c \
  example-code/model-checker/main.c \
  -o /tmp/opencode/model-checker/checker
/tmp/opencode/model-checker/checker \
  example-code/model-checker/stale.graph
\end{verbatim}

Create the output directory first, as shown in the lab README.
The actual failing output is
\begin{verbatim}
unsafe
trace: 0 1 2 3
EX bad: 0x000000000000000c
E[true U bad]: 0x000000000000003f
EG !bad: 0x0000000000000037
\end{verbatim}
Running the same executable on \code{fixed.graph} returns exit status zero and prints
\begin{verbatim}
safe
EX bad: 0x0000000000000008
E[true U bad]: 0x0000000000000008
EG !bad: 0x0000000000000037
\end{verbatim}
Mask \code{0x8} denotes state 3 alone, \code{0x3f} denotes all six states, and \code{0x37} denotes $\{0,1,2,4,5\}$.
The unchanged EG answer is expected: the repair eliminates a bad choice without eliminating the existing infinite safe executions.
These outputs are checked against the API's results by the lab tests.

A compiler integration can export a pass-manager protocol in the same format: states record which analyses are valid, and bad states record a pass consuming invalidated data.
A cache integration can instead generate the six-state abstraction from an explicit transition specification and keep a side table mapping IDs to protocol states and actions.
For either integration, the exporter and abstraction are part of the verification boundary.
A regression gate should preserve the graph, query parameters, and returned trace, and should distinguish all four CLI outcomes.
Treating every nonzero exit code as a discovered compiler bug loses the difference between an invalid model, an incomplete query, and a real model counterexample.

The fixed graph documents a protocol requirement: validity must be checked before entry.
It does not prove that a concrete guard closes a concurrent invalidation race.
If invalidation can happen between the guard and the jump, those events need separate states and interleavings.
If an old code body may be reclaimed while a caller still holds its address, reclamation state must also appear.
If the objective is correctness of reconstructed frames, the abstract state must retain the relevant values or invariants.
Model checking becomes useful when the property, transition granularity, and correspondence to the runtime are reviewed together.

\section{Test with a different algorithm}
\label{sec:ch48-testing}

A checker is particularly vulnerable to tests that merely reproduce its own recurrence.
Comparing one EU fixpoint loop with a differently spelled EU fixpoint loop may preserve the same semantic error.
The lab instead uses finite path enumeration and all-pairs distances as independent oracles.
The test harness runs through every directed graph on one, two, and three labeled states, including all choices of explicit self-loops and missing successors.
There are $2^1+2^4+2^9=530$ such graphs.

For safety, every nonempty initial set is paired with every possible bad-state set.
An independent Floyd--Warshall computation supplies distances between every pair of states.
The oracle takes the minimum distance from an initial state to a bad state, then compares both the status and the returned trace length.
Every unsafe trace is replayed edge by edge.
This gives 28,868 safety queries per test run, including initially bad states, unreachable errors, multiple initial states, cycles, and dead ends.

For EU, a recursive oracle searches simple paths, stopping successfully upon reaching $Q$ and requiring $P$ at all preceding states.
Removing a repeated-state segment preserves an until witness, so this search is complete for the finite graphs under test.
For EG, a separate recursive oracle searches paths with $N$ edges, all of whose states satisfy $P$.
Such a path visits $N+1$ states and must repeat one; its repeated segment can be traversed forever.
Conversely, any infinite $P$ path supplies such a prefix.
EX is checked by scanning individual edges and destination membership.
All operand sets are enumerated, including every ordered pair for EU, giving 41,360 CTL queries per run.

Targeted tests supplement the tiny-graph enumeration.
Twenty malformed streams exercise lexical errors, overflow, missing fields, illegal IDs, duplicate initial IDs, and trailing data.
Additional tests exercise accepted duplicate edges, null arguments, invalid operator values, out-of-universe masks, and zeroed outputs after errors.
A 64-state chain checks bit 63, a 63-edge trace, exact BFS and fixpoint budget boundaries, and \code{SIZE\_MAX} budgets.
A complete 64-state graph is parsed with 64 initial states and all 4096 edge records.
The two cache files are loaded through the real parser, and their safety outcomes and all six printed CTL sets are checked.

The command
\begin{verbatim}
sh example-code/model-checker/test.sh
\end{verbatim}
was run with GCC 16.2.1 and Clang 22.1.8.
Each compiler built the CLI and tests using \code{-std=c11 -Wall -Wextra -Wpedantic -Werror}, first at optimization level \code{-O2}, then with AddressSanitizer and UndefinedBehaviorSanitizer at \code{-O1} with debug information and frame pointers.
All four configurations passed the oracle, boundary, model, and CLI checks; the sanitizer runs reported no errors.
The runner writes executables and captured CLI output under \code{/tmp/opencode/model-checker}.
These are executed tests, not promised future coverage.
They are also finite evidence: the 64-state space was not exhaustively enumerated, stream I/O failures were not fault-injected, and the implementation has not been formally verified.

\section{Limits and useful extensions}
\label{sec:ch48-limits}

The limit of 64 states is checked explicitly rather than hidden behind a shift or an overflowing allocation size.
The universe helper handles $N=64$ separately, because shifting a 64-bit value by 64 is not a way to construct an all-ones mask in C.
The parser bounds stored records, but does not impose a byte budget or a deadline on its input stream.
Query budgets count BFS dequeues and CTL predecessor sweeps; they are deterministic work limits, not wall-clock limits.
These distinctions belong in an API contract because callers may otherwise assume that a bounded graph implies bounded time spent waiting for input.

Larger explicit-state checkers need a representation strategy before they need more temporal syntax.
An adjacency-list graph supports sparse BFS in $O(N+M)$ time.
Reverse edges and a worklist can propagate an EU frontier without rescanning every state each round.
EG can be computed using strongly connected components in the subgraph induced by $P$, retaining states that reach a cycle inside that subgraph.
These are algorithmic alternatives to the small bitset kernel, not claims about optimizations implemented here.

Generating states on demand can avoid constructing unreachable parts of a model for reachability queries.
Hash-consing state descriptions then becomes part of correctness: different concrete model states must not be merged merely because their hashes collide.
Symbolic representations can compress some large sets but do not remove the need for precise transition and proposition semantics.
Partial-order reduction requires conditions connecting the reduced interleavings to the properties being checked.
Fairness adds obligations beyond accepting any reachable cycle.
Each extension changes the verification argument, the output contract, or both; a larger checker should preserve the small lab as a reference oracle on bounded instances.

\section{Exercises}
\label{sec:ch48-exercises}

Answers are in \code{example-code/model-checker/SOLUTIONS.md}.
\begin{enumerate}
  \item \label{ex:ch48-dead-end}
    A graph contains only initial state 0 and no explicit edge.
    Compute $EX\,p$, $EG\,p$, and $E[\mathit{true}\ U\ p]$ when $p$ holds at 0 and when it does not.
    Explain how an exporter could preserve the ability to detect deadlock after totalization.
  \item \label{ex:ch48-until-rounds}
    Compute every EU approximation for the buggy cache graph with $P$ equal to all states and $Q=\{3\}$.
    How many predecessor sweeps does the implementation need to return success?
    What does a budget of four return, and what changes for the repaired graph?
  \item \label{ex:ch48-existential-safety}
    Exhibit an infinite safe execution of the buggy graph and a finite failing execution from the same initial state.
    Give a set expression using EU that checks universal safety of all initial states.
  \item \label{ex:ch48-multiple-roots}
    Use four states with initial set $\{0,2\}$, edges $0\to1$, $1\to3$, and $2\to3$, and bad set $\{3\}$.
    Work through the queue and parents, and give the shortest returned trace.
    Explain why overwriting a parent on a second discovery would break the argument.
  \item \label{ex:ch48-capacity}
    For the chain $0\to1\to\cdots\to63$ with only state 63 bad, determine the smallest budgets that let BFS and $EU(\mathit{true},\mathit{bad})$ complete.
    Repeat for $EG(\{0,\ldots,62\})$.
    Justify the fixed queue and trace capacities without appealing to a successful test run.
  \item \label{ex:ch48-refinement}
    Refine the cache model so that checking an entry guard and entering code are separate transitions.
    Describe an invalidation interleaving omitted by the coarse model.
    State what additional runtime protocol must be modeled before claiming that the repair handles that interleaving.
\end{enumerate}

\section{Further reading}
\label{sec:ch48-reading}

Baier and Katoen~\cite{baier-2008-model-checking} develop transition-system semantics, temporal logics, and the algorithms that connect formulas to finite graphs.
Huth and Ryan~\cite{huth-2004-logic} provide a useful bridge from logical specifications to practical model-checking computations and their correctness arguments.
The \emph{Handbook of Model Checking}~\cite{clarke-2018-handbook} places explicit-state checking alongside symbolic techniques, abstraction, reduction, and other approaches needed when a model outgrows this lab's deliberately small representation.
